\section{变分态制备}
\label{sec:ansatz}

\subsection{示意起点与终点}
\label{sec:points}

三对示意起点/终点，每相一对，标在 $\tilde{Z}_{\mathcal R}$ 相图上，
见图~\ref{fig:D05}：平庸相 $(0.01,0.02)\to(0.21,0.42)$，
$\tilde{Z}_{\mathcal R}\approx0.996$；拓扑相 $(0.99,0.02)\to(0.90,0.20)$，
$\tilde{Z}_{\mathcal R}\approx-0.98$；AFM 相 $(0.50,1.98)\to(0.61,1.44)$，
$\tilde{Z}_{\mathcal R}\approx0.009$。
终点是相深处的经典参考目标，待相图落定后选取——变分运行的示范终点，
并非真实演化终点本身。终点置于相深处，使参考路径避开有限尺寸漂移最大的
边界区。起点沿用三色圆点样式；终点标记形状颜色均区别于起点，
箭头给出每对的参考方向。

\begin{figure}[htbp]
\centering
\includegraphics[width=\columnwidth]{exp01_D05.pdf}
\caption{图~\ref{fig:D01} $\tilde{Z}_{\mathcal R}$ 相图上的示意起点/终点对，
每相一对、参考终点在相深处；颜色区分各对，箭头标参考方向。
\label{fig:D05}}
\end{figure}

\subsection{变分电路结构}

我们的变分电路遵循变分本征求解器范式~\cite{peruzzo2014variational}：
制备初态，再作用参数化拟设。经典循环最小化
\begin{equation}
E(\theta)=\langle\psi(\theta)|H|\psi(\theta)\rangle,
\end{equation}
其中 $|\psi(\theta)\rangle$ 为参数 $\theta$ 下的电路输出。
这类变分算法见综述~\cite{cerezo2021variational,tilly2022variational}，
其中收集了代价函数设计、拟设族与最佳实践。
实际瓶颈是优化景观而非表达能力：随机电路形成 $2$-design 时梯度随比特数
指数消失~\cite{mcclean2018barren}；硬件高效电路随深度出现陡峭的可训练性
转变~\cite{campos2021abrupt}；最优参数可以问题规模的逆多项式集中~\cite{akshay2021parameter}；
光滑的哈密顿量变分解跨尺寸迁移且梯度不消失~\cite{mele2022avoiding}；
贫瘠高原区还可藏指数多个陷阱~\cite{nemkov2025barren}。
Krylov 型迭代求解器是到基态的另一条路~\cite{bharti2021iterative}。
硬件高效变分优化最早在超导真机上对小分子与量子磁体实现~\cite{kandala2017hardware}，
自旋链变分模拟随后在云超导设备上推到 $102$ 比特~\cite{yu2023simulating}。
我们的 orbit 拟设与之都不同：以反射对称性绑定参数、初态定纠缠结构，详述如下。
三初态（奇键单态乘积及其对应态，均在 $P=-2$ 扇区，$P$ 标记组合对称性本征值）
配 orbit 拟设，镜面对映键对共享转角（奇键 $O_j\leftrightarrow O_{M+1-j}$，
偶键 $E_j\leftrightarrow E_{M-j}$）。即拟设分两步构造：先以相匹配初态定下
纠缠结构，再给所有键穿上由反射对称性绑定的转角，参数量随键 orbit 数而非键数增长。
每相一个示意图见图~\ref{fig:D06}(a)--(c)：初态 + 所需拟设电路，分别为平庸、
拓扑、AFM 相。
该构造保证试探态永不离开目标对称性扇区；交叠是否足够，由第~\ref{sec:comparison}~节的
真机-模拟机基准裁决。

\begin{figure}[htbp]
\centering
\includegraphics[width=\columnwidth]{exp04_D06a.pdf}
\caption{(a) 平庸初态 + 所需拟设电路。虚线框 $U^{(k)}$ 为第 $k$ 变分层；
$\theta^{(k)}_j$（$\phi^{(k)}_j$）为第 $k$ 层第 $j$ 镜面对映键 orbit 共用的
$XX{+}YY$（$ZZ$）转角（奇键 $O_j\leftrightarrow O_{M+1-j}$，偶键
$E_j\leftrightarrow E_{M-j}$；$\delta=0$ 时 $\theta=\phi$）。偶键子层先作用。
\label{fig:D06}}
\end{figure}

\begin{figure}[htbp]
\centering
\addtocounter{figure}{-1}
\includegraphics[width=\columnwidth]{exp04_D06b.pdf}
\caption{(b) 拓扑初态 + 所需拟设电路，记号同 (a)；奇键子层先作用。}
\end{figure}

\begin{figure}[htbp]
\centering
\addtocounter{figure}{-1}
\includegraphics[width=\columnwidth]{exp04_D06c.pdf}
\caption{(c) AFM 初态 + 所需拟设电路，记号同 (a)；偶键子层先作用。}
\end{figure}
